% Table source-order #20: Confidence-triggered cascade. Confidence-triggered cascade, vision → s
\begin{table}[tp]
\centering
\caption{Confidence-triggered cascade. Confidence-triggered cascade, vision → som. The escalation decision sees only the cheap run's own episode: no outcome, no rich-run information. In 0 of the 6 comparable cells does any operating point beat always-rich on success rate; the cells marked \emph{rich is worse} are excluded because the cascade's premise fails there, and that exclusion rule (\texttt{cheap\_sr\ \textgreater{}=\ rich\_sr}) is outcome-dependent. \textbf{Caution:} The two counts are split for a reason. A point also `wins' by matching always-rich's SR at lower cost, and every such point here is a tie artefact: on \texttt{WA·B1} all of them sit at exactly the rich arm's SR while 44-60 of 104 episodes are tied at the cutoff, so the ranking falls through to task id and the reachable SR spans 8.65-14.42. Counting the two together is what let this table print a WA win for months while the prose said the WA exception was withdrawn: the rows disagreed with the sentence and nobody could see it, because the WA rows were dropped before render until 2026-08-03. \textbf{Caution:} Every number is an offline splice: an escalated task takes its outcome from a standalone rich run, whereas a real cascade would start the rich episode after the cheap one had already acted on a stateful site. That sequential outcome is unobserved in this project. Source: \texttt{confidence\_cascade\_with\_wa.json}}
\label{tab:t20}
\footnotesize
\setlength{\tabcolsep}{3pt}
\begin{tabular}{@{}
  >{\raggedright\arraybackslash}p{(\linewidth - 16\tabcolsep) * \real{0.1111}}
  >{\raggedright\arraybackslash}p{(\linewidth - 16\tabcolsep) * \real{0.1111}}
  >{\raggedright\arraybackslash}p{(\linewidth - 16\tabcolsep) * \real{0.1111}}
  >{\raggedright\arraybackslash}p{(\linewidth - 16\tabcolsep) * \real{0.1111}}
  >{\raggedright\arraybackslash}p{(\linewidth - 16\tabcolsep) * \real{0.1111}}
  >{\raggedright\arraybackslash}p{(\linewidth - 16\tabcolsep) * \real{0.1111}}
  >{\raggedright\arraybackslash}p{(\linewidth - 16\tabcolsep) * \real{0.1111}}
  >{\raggedright\arraybackslash}p{(\linewidth - 16\tabcolsep) * \real{0.1111}}
  >{\raggedright\arraybackslash}p{(\linewidth - 16\tabcolsep) * \real{0.1111}}@{}}
\toprule
\begin{minipage}[b]{\linewidth}\raggedright
cell
\end{minipage} & \begin{minipage}[b]{\linewidth}\raggedright
n
\end{minipage} & \begin{minipage}[b]{\linewidth}\raggedright
cheap SR
\end{minipage} & \begin{minipage}[b]{\linewidth}\raggedright
rich SR
\end{minipage} & \begin{minipage}[b]{\linewidth}\raggedright
always-rich cost
\end{minipage} & \begin{minipage}[b]{\linewidth}\raggedright
oracle SR
\end{minipage} & \begin{minipage}[b]{\linewidth}\raggedright
comparable?
\end{minipage} & \begin{minipage}[b]{\linewidth}\raggedright
beats always-rich --- strictly / SR-tied
\end{minipage} & \begin{minipage}[b]{\linewidth}\raggedright
signals dropped
\end{minipage} \\
\midrule
cls·B0 & 224 & 25.00 & 27.23 & 1.116× & 33.93 & yes & 0 / 0 & 2 \\
cls·B1 & 224 & 12.50 & 14.29 & 1.397× & 19.20 & yes & 0 / 0 & 0 \\
cls·B2 & 224 & 2.23 & 2.23 & 1.284× & 4.46 & --- \emph{(rich is worse)} & \textbf{11} / 22 & 0 \\
red·B0 & 203 & 7.39 & 14.78 & 1.126× & 18.23 & yes & 0 / 0 & 0 \\
red·B1 & 203 & 2.46 & 7.39 & 1.527× & 8.37 & yes & 0 / 0 & 2 \\
red·B2 & 203 & 1.97 & 0.99 & 1.633× & 2.96 & --- \emph{(rich is worse)} & \textbf{41} / 5 & 1 \\
WA·B0 & 104 & 19.23 & 22.12 & 1.054× & 29.81 & yes & 0 / 0 & 0 \\
WA·B1 & 104 & 9.62 & 13.46 & 1.778× & 16.35 & yes & 0 / 3 & 1 \\
\bottomrule
\end{tabular}
\end{table}
